\chapter{Martyngały ze skokami, kompensacja i granica zakresu tomu II}
\label{ch:jump-mart}

W całym rozdziale przyjmujemy warunki zwykłe filtracji
(rozdział~\ref{ch:mart-cont}). Procesy c\`adl\`ag
rozpatrujemy w wersjach o takich trajektoriach
poza jednym zbiorem zerowym; na tym zbiorze
można przyjąć zerową trajektorię dzięki zupełności.
Konstrukcję ograniczamy do przypadków, w których
można uzasadnić wszystkie własności bez ogólnej
teorii całkowania względem półmartyngałów.

\section{Dlaczego trzeba oddzielić dwa rodzaje wariacji kwadratowej?}
Dla ciągłego lokalnego martyngału jego wariacja kwadratowa
jest ciągłym procesem rosnącym i opisuje akumulację
losowych fluktuacji. Przy skokach pojawia się dodatkowe
zjawisko: pojedynczy skok o wysokości $\Delta X_t$
wnosi do wariacji kwadratowej dokładnie $(\Delta X_t)^2$.
Dla martyngałów z częścią skokową trzeba odróżniać
rzeczywistą wariację kwadratową od jej przewidywalnej
kompensacji. Nie jest ona na ogół warunkową średnią
całego nawiasu w jednej ustalonej chwili. Obie wielkości są użyteczne, ale w ogólności
nie są sobie równe.

\begin{definicja}[Skok i suma kwadratów skoków]
Dla procesu c\`adl\`ag $X$ definiujemy skok
$\Delta X_t=X_t-X_{t-}$, $t>0$. Dla trajektorii lokalnie skończonej wariacji
całkowitej poniżej udowodnimy, że jej wariacja
kwadratowa wynosi
\[
[X]_t=\sum_{0<s\le t}(\Delta X_s)^2.
\]
Dla szerszych klas procesów granica sum kwadratów
przyrostów wymaga dodatkowej konstrukcji, stanowiącej
część teorii półmartyngałów rozwijanej w tomie III.
\end{definicja}


\begin{definicja}[Wariacja kwadratowa w klasie procesów z niniejszego rozdziału]
Dla rzeczywistego procesu c\`adl\`ag $X$ i podziału
$\pi=\{0=t_0<\cdots<t_m=T\}$ definiujemy
$V_t^\pi(X)=\sum_i(X_{t\wedge t_{i+1}}-X_{t\wedge t_i})^2$.
Jeśli dla każdego ciągu deterministycznych podziałów o
średnicach dążących do zera procesy $V^\pi(X)$ zbiegają
jednostajnie na $[0,T]$ według prawdopodobieństwa do
procesu c\`adl\`ag $A$, piszemy $[X]=A$ na $[0,T]$.
Jednoznaczność z dokładnością do nierozróżnialności
wynika z nierówności trójkąta dla supremum i
przeliczalnej rodziny horyzontów, jak w rozdziale
\ref{ch:quadratic-cont}.
Istnienie granicy wykażemy dla procesów o lokalnie
skończonej wariacji i ich sum z ciągłymi procesami
o istniejącej wariacji kwadratowej. Ogólną teorię
półmartyngałów pozostawiamy tomowi III.
Ważne jest użycie niedokończonego przyrostu:
gdy $t$ jest chwilą skoku, jego kwadrat jest
uwzględniony już w chwili $t$, a nie dopiero
w następnym węźle podziału.
\end{definicja}

\begin{lemat}[Wariacja kwadratowa funkcji o skończonej wariacji]
\label{lem:cadlag-finite-variation-qv}
Niech $f:[0,T]\to\R$ będzie funkcją c\`adl\`ag
o skończonej wariacji całkowitej.
Dopuszczamy nieskończenie wiele skoków. Dla dowolnych deterministycznych podziałów
o średnicach malejących do zera zachodzi
\[
\sup_{t\le T}\left|V_t^\pi(f)-
\sum_{0<s\le t}(\Delta f_s)^2\right|\longrightarrow0.
\]
W szczególności zbieżność zachodzi również w chwilach
skoku, z prawostronną konwencją $f_t$ jako wartości
trajektorii już po skoku.
\end{lemat}
\begin{proof}
Najpierw zakładamy skończoną liczbę skoków.
Rozłóżmy $f=g+j$, gdzie
$j(t)=\sum_{0<s\le t}\Delta f_s$,
a $g=f-j$ jest ciągła i ma skończoną wariację
całkowitą. Dla dostatecznie drobnych podziałów
każda półotwarta komórka $(t_i,t_{i+1}]$
zawiera co najwyżej jeden skok.
Suma $(\Delta j)^2$ po pełnych komórkach oraz
aktualnej komórce jest dokładnie sumą kwadratów
skoków, które wystąpiły nie później niż $t$. Suma $(\Delta g)^2$ dąży do zera
z lematu~\ref{lem:finite-var-zero-quadratic}. Wyrazy mieszane są
ograniczone przez nierówność Cauchy'ego--Schwarza:
\[
\left|2\sum\Delta g\Delta j\right|
\le2\Bigl(\sum(\Delta g)^2\Bigr)^{1/2}
\Bigl(\sum(\Delta j)^2\Bigr)^{1/2}\longrightarrow0.
\]
W chwili skoku $t$ wartość $f_t$ zawiera już skok,
a więc odpowiadający mu kwadrat pojawia się również
w sumie $V_t^\pi$. Oszacowania są jednostajne w $t$,
ponieważ liczba skoków jest skończona, a funkcja $g$
jest jednostajnie ciągła.

W ogólnym przypadku dla każdego skończonego zbioru
chwil skoku suma modułów skoków jest nie większa
od $\operatorname{Var}_{[0,T]}(f)$: dobieramy rozłączne
przedziały kończące się w tych chwilach i przechodzimy
z lewymi końcami do chwil skoków, używając lewych granic.
W szczególności dla każdego $\varepsilon>0$ tylko
skończenie wiele skoków ma moduł większy od
$\varepsilon$, wszystkie niezerowe skoki tworzą
zbiór przeliczalny oraz
$\sum_{0<s\le T}|\Delta f_s|\le\operatorname{Var}_{[0,T]}(f)$.
Szereg $j(t)=\sum_{0<s\le t}\Delta f_s$
jest jednostajnie zbieżny i określa funkcję
c\`adl\`ag o skończonej wariacji z tymi samymi
skokami co $f$. Istotnie sumy skończone mają
te własności, a jednostajna granica zachowuje
prawostronną ciągłość i lewe granice.
Wariacja ogona jest nie większa niż suma modułów
odrzuconych skoków. Funkcja $g=f-j$ jest ciągła.
Wybierzmy skończone sumy skokowe $j_n$ z
$\operatorname{Var}(j-j_n)\to0$.
Dla każdego podziału i $t\le T$ nierówność
trójkąta dla normy euklidesowej daje
\[
\left|\sqrt{V_t^\pi(f)}-\sqrt{V_t^\pi(g+j_n)}\right|
\le\sqrt{V_t^\pi(j-j_n)}
\le\operatorname{Var}_{[0,T]}(j-j_n).
\]
Obie pierwiastkowane sumy są jednostajnie
ograniczone przez
$\operatorname{Var}(g)+\operatorname{Var}(j)+
\sum_{0<s\le T}|\Delta f_s|$.
Zatem także różnica samych kwadratów jest
jednostajnie mała przy dużym $n$, niezależnie
od siatki i $t$. Suma kwadratów skoków
$j_n$ zbiega jednostajnie do sumy kwadratów
wszystkich skoków, gdyż szereg kwadratów
jest zbieżny. Stosujemy przypadek skończony
przy ustalonym $n$, a następnie przechodzimy
z $n$ do nieskończoności.
\end{proof}

\begin{definicja}[Kompensator i przewidywalny nawias]
\label{def:jump-predictable-bracket}
Dla adaptowanego niemalejącego procesu c\`adl\`ag
$V$ z $V_0=0$ kompensatorem nazywamy przewidywalny,
niemalejący proces c\`adl\`ag $A$ z $A_0=0$,
taki że $V-A$ jest lokalnym martyngałem.
W przypadku ciągłych lokalnych martyngałów
$[M]$ jest już przewidywalny i sam jest swoim
kompensatorem. Dla lokalnie kwadratowo
całkowalnego martyngału $M$ z $M_0=0$
przewidywalny nawias $\langle M\rangle$,
gdy istnieje, jest przewidywalnym niemalejącym
procesem z wartością początkową zero, takim że
$M^2-\langle M\rangle$ jest lokalnym martyngałem.
W rozważanych niżej przykładach sprawdzimy
bezpośrednio, że również $M^2-[M]$ jest
lokalnym martyngałem, a zatem $\langle M\rangle$
jest kompensatorem $[M]$.
Jednoznaczność udowodnimy poniżej. Ogólne
twierdzenie o istnieniu kompensatora wymaga
dodatkowej teorii (twierdzenia Dooba--Meyera)
i nie jest tutaj zakładane.
\end{definicja}

\begin{twierdzenie}[Nawias i kompensator procesu Poissona]
\label{thm:poisson-brackets}
Niech $N$ będzie procesem Poissona o intensywności
$\lambda>0$, w jego uzupełnionej filtracji naturalnej,
a $M_t=N_t-\lambda t$. Wtedy $M$ jest kwadratowo
całkowalnym martyngałem z
\[
[M]_t=N_t,\qquad \langle M\rangle_t=\lambda t.
\]
Pierwszy proces jest wariacją kwadratową trajektorii,
a drugi \emph{przewidywalnym nawiasem}: jest ciągły,
niemalejący i przewidywalny, a $M_t^2-\langle M\rangle_t$
jest martyngałem. Jednoznaczność tej kompensacji w klasie procesów
przewidywalnych wynika z wniosku
\ref{cor:predictable-compensator-unique}.
Ponadto $\E M_t^2=\lambda t$.
\end{twierdzenie}
\begin{proof}
Proces $M$ jest kwadratowo całkowalnym martyngałem,
gdyż przyrosty $N$ są niezależne od przeszłości
i mają średnią $\lambda(t-s)$ oraz skończony
drugi moment (rozdział~\ref{ch:poisson-process}).
Trajektorie $M$ mają skoki równe $1$ w chwilach
zdarzeń Poissona i liniowe odcinki o nachyleniu
$-\lambda$ pomiędzy nimi. Na każdym skończonym
przedziale liczba skoków jest skończona prawie na pewno,
więc poprzedni lemat daje
$[M]_t=\sum_{0<s\le t}1=N_t$.
Dla $s<t$, pisząc $U=N_t-N_s$, mamy
$M_t=M_s+(U-\lambda(t-s))$.
Zmienna $U$ jest niezależna od $\calF_s$, a jej
średnia i wariancja wynoszą $\lambda(t-s)$.
Stąd
\[
\E[M_t^2\mid\calF_s]
=M_s^2+\lambda(t-s),
\]
co pokazuje, że $M_t^2-\lambda t$ jest martyngałem.
Deterministyczna funkcja $t\mapsto\lambda t$ jest
przewidywalna. Równanie przy $s=0$ daje też
$\E M_t^2=\lambda t$.
Ponieważ $[M]-\lambda t=N-\lambda t=M$
jest martyngałem, także $M^2-[M]$
jest martyngałem jako różnica dwóch martyngałów.
\end{proof}


\begin{lemat}[Jednoznaczność ciągłego kompensatora kwadratu]
\label{lem:continuous-compensator-unique}
Niech $M$ będzie lokalnym martyngałem c\`adl\`ag, a $A$ i $B$
ciągłymi, adaptowanymi procesami o lokalnie skończonej wariacji,
$A_0=B_0=0$. Jeżeli zarówno $M^2-A$, jak i $M^2-B$
są lokalnymi martyngałami, to $A$ i $B$ są nierozróżnialne.
\end{lemat}
\begin{proof}
Różnica $D=A-B=(M^2-B)-(M^2-A)$ jest ciągłym
lokalnym martyngałem o skończonej wariacji na przedziałach
zwartych. Z twierdzenia~\ref{thm:continuous-local-finite-variation}
wynika, że każdy taki proces jest stały. Ponieważ $D_0=0$,
proces $D$ jest nierozróżnialny od zera.
\end{proof}

\begin{lemat}[Całkowanie względem martyngału o całkowalnej wariacji]
\label{lem:predictable-fv-mean-zero}
Niech $D$ będzie rzeczywistym martyngałem c\`adl\`ag
z $D_0=0$ oraz $\E\operatorname{Var}_{[0,T]}(D)<\infty$.
Dla każdego ograniczonego procesu przewidywalnego $H$
zachodzi
\[
\E\int_{(0,T]} H_t\,\mathrm dD_t=0.
\]
Całka jest tutaj zwykłą trajektoryjną całką
względem miary znakowanej Lebesgue'a--Stieltjesa,
a nie całką It\^o.
\end{lemat}
\begin{proof}
Dla generatorów $\chi_F(\omega)\chi_{(a,b]}(t)$,
$F\in\calF_a$ i $0\le a<b\le T$, lewa strona
wynosi $\E[\chi_F(D_b-D_a)]=0$.
Rodzina tych prostokątów, wraz ze zbiorem pustym,
jest zamknięta na skończone przecięcia i generuje
$\sigma$-ciało przewidywalne ograniczone do
$(0,T]\times\Omega$ (rozdział~\ref{ch:filtracje}).
Dla każdego zbioru z tego $\sigma$-ciała całka
jego funkcji charakterystycznej jest mierzalna:
najpierw dla prostokątów, a następnie przez
twierdzenie o klasie monotonicznej i trajektoryjną
zbieżność zmajoryzowaną względem miary wariacji.
Funkcja zbioru otrzymana przez wzięcie wartości
oczekiwanej jest skończoną miarą znakowaną,
ponieważ całkowalna wariacja majoruje moduły
całek i pozwala przechodzić do granicy dla
sum rozłącznych. Jej wartość na całej przestrzeni
wynosi $\E D_T=0$. Z jednoznaczności miar
na układzie generującym (lub twierdzenia
$\pi$--$\lambda$ dla zbiorów, na których
miara ma wartość zero) jest identycznie zerowa.
Aproksymacja ograniczonego $H$ funkcjami
prostymi kończy dowód.
\end{proof}

\begin{twierdzenie}[Przewidywalny lokalny martyngał o skończonej wariacji]
\label{thm:predictable-fv-local-constant}
Jeżeli $D$ jest przewidywalnym lokalnym martyngałem
c\`adl\`ag o lokalnie skończonej wariacji całkowitej,
to $D$ jest nierozróżnialny od procesu stałego $D_0$.
\end{twierdzenie}
\begin{proof}
Odejmujemy $D_0$ i niech
$V_t=\operatorname{Var}_{[0,t]}(D)$.
Proces $V$ jest adaptowany, niemalejący i c\`adl\`ag:
wariację można obliczać po podziałach z wymiernymi
punktami wewnętrznymi oraz końcem $t$ dzięki
prawostronnej ciągłości; miara wariacji trajektorii
ma atom w $t>0$ o masie $|\Delta D_t|$.
Zatem $\Delta V_t=|\Delta D_t|$.
Niech $(\sigma_n)$ będzie ciągiem lokalizującym $D$
i połóżmy
\[
\tau_n=\sigma_n\wedge n\wedge
\inf\{t\ge0: V_t\ge n\}.
\]
Są to czasy zatrzymania, gdyż dla niemalejącego
prawostronnie ciągłego $V$ zdarzenie trafienia
do poziomu $n$ do chwili $t$ jest równe
$\{V_t\ge n\}$. Mamy $\tau_n\uparrow\infty$ p.n.
Proces $D^{\tau_n}$ jest martyngałem na mocy
twierdzenia~\ref{thm:mart-cont-stopping}
zastosowanego do $D^{\sigma_n}$.
Ponadto jest przewidywalny: dla dowolnego
czasu zatrzymania $\tau$ rozkład
\[
D_{t\wedge\tau}=D_t\chi_{\{t\le\tau\}}+
D_\tau\chi_{\{\tau<t\}}
\]
składa się z procesów przewidywalnych
(w drugim składniku przyjmujemy zero przed
$\tau$, także gdy $\tau=\infty$).
Wskaźnik pierwszego zdarzenia jest lewostronnie
ciągły i adaptowany, podobnie drugi składnik
jest lewostronnie ciągłym procesem adaptowanym.
Wartość w zerze jest mierzalna względem $\calF_0$.

Przed $\tau_n$ wariacja jest ograniczona przez $n$,
więc $V_{\tau_n}\le n+|\Delta D_{\tau_n}|$
(przy $\tau_n=0$ przyjmujemy skok równy zero).
Ponieważ $|D_{\tau_n-}|\le n$ i $D_{\tau_n}$
jest całkowalna na mocy ograniczonego próbkowania
opcjonalnego dla $D^{\sigma_n}$, otrzymujemy
$\E V_{\tau_n}<\infty$.
Zatrzymany proces ma zatem całkowalną wariację
na każdym horyzoncie.
W dalszej części oznaczmy go przez $L$.

Proces $L_-$ jest przewidywalny jako adaptowany
proces o lewostronnie ciągłych trajektoriach
dla $t>0$. Stąd
$H_t=\operatorname{sgn}(\Delta L_t)$
jest ograniczony i przewidywalny
(przyjmujemy $H_0=0$ oraz $\operatorname{sgn}0=0$).
Każda trajektoria o skończonej wariacji
ma co najwyżej przeliczalnie wiele skoków.
Ponieważ $H$ znika poza nimi, mamy
\[
\int_{(0,T]}H_t\,\mathrm dL_t
=\sum_{0<t\le T}|\Delta L_t|.
\]
Lemat~\ref{lem:predictable-fv-mean-zero}
daje średnią zero tej nieujemnej sumy,
więc $L$ nie ma skoków p.n. na $[0,T]$.
Po przecięciu zdarzeń pełnej miary dla
całkowitych $T$ proces $L$ jest ciągły.
Twierdzenie~\ref{thm:continuous-local-finite-variation}
daje $L=0$ z dokładnością do nierozróżnialności.
Przeliczalność lokalizacji i $\tau_n\uparrow\infty$
kończą dowód dla $D$.
\end{proof}

\begin{wniosek}[Jednoznaczność kompensatorów przewidywalnych]
\label{cor:predictable-compensator-unique}
Jeśli $A,B$ są kompensatorami tego samego
procesu $V$ z definicji
\ref{def:jump-predictable-bracket}, to są
nierozróżnialne. Także przewidywalny nawias,
jeśli istnieje, jest jednoznaczny.
\end{wniosek}
\begin{proof}
Różnica $A-B$ jest przewidywalnym lokalnym
martyngałem o lokalnie skończonej wariacji
z wartością początkową zero, zarówno gdy
odejmujemy dwa kompensatory, jak i dwa nawiasy
przewidywalne. Stosujemy twierdzenie
\ref{thm:predictable-fv-local-constant}.
\end{proof}

\begin{uwaga}[Co oznacza kompensator?]
Wzór $\langle M\rangle_t=\lambda t$ nie oznacza,
że losowa liczba skoków $N_t$ jest równa $\lambda t$.
Oznacza, że $[M]_t-\langle M\rangle_t=N_t-\lambda t$
jest martyngałem: po odjęciu przewidywalnego tempa
wzrostu pozostaje wycentrowana fluktuacja.
Przewidywalne kompensatory mogą też mieć skoki.
Ich jednoznaczność wynika z wniosku
\ref{cor:predictable-compensator-unique};
ogólną konstrukcję pozostawiamy tomowi III.
\end{uwaga}


\begin{uwaga}[Różnica między nawiasami jest rzeczywista]
W chwili skoku $N$ wariacja kwadratowa $[M]$ skacze
o jeden, a przewidywalny nawias $\langle M\rangle$
pozostaje ciągły. Różnica $N_t-\lambda t=M_t$
sama jest martyngałem. Nie należy zatem utożsamiać
$[M]$ i $\langle M\rangle$, chociaż ich wartości
oczekiwane są równe: $\E\bigl([M]_t\bigr)=\E\bigl(\langle M\rangle_t\bigr)=\lambda t$.
\end{uwaga}

\begin{przyklad}[Przewidywalny skok nawiasu]
\label{ex:predictable-bracket-jump}
Niech $Y$ przyjmuje wartości $-1,1$ z równymi
prawdopodobieństwami, a filtracja będzie trywialna
przed chwilą $1$ i równa $\sigma(Y)$ od chwili $1$
(z uzupełnieniem). Jest to filtracja prawostronnie
ciągła. Proces $M_t=Y\chi_{[1,\infty)}(t)$
jest martyngałem: przed ujawnieniem $Y$ jego
warunkowa średnia jest zerowa, a potem proces
jest stały.
W tym przypadku
\[
[M]_t=Y^2\chi_{[1,\infty)}(t)
=\chi_{[1,\infty)}(t)=\langle M\rangle_t.
\]
Proces deterministyczny $\chi_{[1,\infty)}$
jest przewidywalny, gdyż zbiór $[1,\infty)$
jest przeliczalnym przecięciem przedziałów
$(1-1/n,\infty)$, mierzalnych w
$\sigma$-ciele przewidywalnym (rozdział
\ref{ch:filtracje}). Różnica
$M^2-\langle M\rangle$ jest zerowa.
Nawias przewidywalny może więc mieć skok,
a procesy o skokach nie muszą mieć różnych
nawiasów. W tym przykładzie nie istnieje ciągły
kompensator kwadratu nawet w sensie lokalnym:
każdy ciągły adaptowany kandydat jest
przewidywalny i z jednoznaczności musiałby
być nierozróżnialny od $\chi_{[1,\infty)}$.
\end{przyklad}

\begin{center}
\begin{tikzpicture}[x=1.5cm,y=0.8cm,>=Latex]
\draw[->] (0,0)--(4.5,0) node[right] {$t$};
\draw[->] (0,0)--(0,4.5) node[above] {nawias};
\draw[blue!75!black,very thick] (0,0)--(.7,0) (.7,1)--(1.6,1) (1.6,2)--(3.2,2) (3.2,3)--(4,3);
\foreach \x/\y in {.7/1,1.6/2,3.2/3}{\fill[blue!75!black] (\x,\y) circle (1.7pt); \draw[blue!75!black,fill=white] (\x,{\y-1}) circle (1.7pt);}
\draw[orange!80!black,thick,dashed] (0,0)--(4,4);
\node[blue!75!black,below] at (3.7,3) {$[M]_t=N_t$};
\node[orange!80!black,above left] at (3.6,3.6) {$\langle M\rangle_t=t$};
\end{tikzpicture}
\par\smallskip\small Przykładowa trajektoria nawiasu
kompensowanego procesu Poissona o intensywności $1$.
Nawias przewidywalny jest deterministyczny, choć
zrealizowana liczba skoków jest losowa.
\end{center}

\section{Czasy przewidywalne i skoki procesów zliczających}

\begin{definicja}[Czas przewidywalny i całkowicie nieosiągalny]
Czas zatrzymania $\tau$ nazywamy przewidywalnym,
jeżeli istnieje niemalejący ciąg czasów zatrzymania
$\tau_n\uparrow\tau$ p.n. oraz $\tau_n<\tau$
na $\{\tau>0\}$ p.n. (dla $\tau=\infty$
żądamy $\tau_n<\infty$).
Czasy $\tau_n$ nazywamy ciągiem zapowiadającym.
Czas zatrzymania $S$ nazywamy całkowicie
nieosiągalnym, jeżeli dla każdego
przewidywalnego czasu zatrzymania $\tau$
\[
\Prob(0<S=\tau<\infty)=0
\]
oraz $\Prob(S=0)=0$.
Dodatnia deterministyczna chwila jest
przewidywalna, np. przez ciąg $(t-1/n)^+$.
\end{definicja}

\begin{lemat}[Wykres czasu przewidywalnego]
\label{lem:predictable-time-graph}
Jeżeli $\tau$ jest przewidywalnym czasem zatrzymania,
to proces $\chi_{\{0<t=\tau<\infty\}}$
jest przewidywalny.
\end{lemat}
\begin{proof}
Dla dowolnych czasów zatrzymania $\alpha,\beta$
proces $\chi_{\{\alpha<t\le\beta\}}$
jest lewostronnie ciągły i adaptowany
na $(0,\infty)$, zatem przewidywalny.
Dla ciągu zapowiadającego mamy punktowo
poza wspólnym zbiorem zerowym
\[
\{0<t=\tau<\infty\}
=\bigcap_{n\ge1}\{\tau_n<t\le\tau\},
\qquad t>0.
\]
Przy $\tau=\infty$ żaden skończony $t$
nie należy do wszystkich zbiorów po prawej.
Zupełność filtracji umożliwia zmianę
na wspólnym zbiorze zerowym.
\end{proof}

\begin{twierdzenie}[Tożsamość kompensacyjna]
\label{thm:counting-compensation}
Niech $N$ będzie adaptowanym procesem zliczającym
z $N_0=0$ i jednostkowymi skokami, nieeksplodującym,
a $A$ przewidywalnym niemalejącym procesem
c\`adl\`ag z $A_0=0$. Załóżmy, że $N-A$
jest martyngałem oraz $\E N_T+\E A_T<\infty$
dla każdego $T<\infty$. Wtedy dla każdego
nieujemnego procesu przewidywalnego $H$
\[
\E\sum_{0<s\le T}H_s\Delta N_s
=\E\int_{(0,T]}H_s\,\mathrm dA_s
\]
z dopuszczeniem wartości nieskończonej.
Dla procesu Poissona o intensywności $\lambda$
w jego uzupełnionej filtracji naturalnej
prawa strona jest równa
$\lambda\E\int_0^T H_s\,\mathrm d\lambda_1(s)$.
Jeśli $A$ jest ciągły, wszystkie kolejne czasy
skoków $N$ są całkowicie nieosiągalne.
\end{twierdzenie}
\begin{proof}
Dla ograniczonego $H$ stosujemy lemat
\ref{lem:predictable-fv-mean-zero} do $D=N-A$,
którego wariacja na $[0,T]$ jest
ograniczona przez $N_T+A_T$.
Całka względem $N$ jest sumą po chwilach
jego skoków. Dla dowolnego nieujemnego $H$
przechodzimy od $H\wedge n$ przez zbieżność
monotoniczną po obu stronach.
Przypadek Poissona wynika z twierdzenia
\ref{thm:poisson-brackets}.

Dla przewidywalnego $\tau$ stosujemy tożsamość
do procesu $H_t=\chi_{\{0<t=\tau<\infty\}}$,
przewidywalnego z lematu
\ref{lem:predictable-time-graph}.
Przy ciągłym $A$ jego miara nie ma atomów,
więc prawa strona wynosi zero.
Otrzymujemy $\Prob(0<\tau\le T,\Delta N_\tau=1)=0$.
Przechodząc po całkowitych $T$ do nieskończoności,
widzimy, że żaden kolejny czas skoku
$S_k=\inf\{t:N_t\ge k\}$ nie może z dodatnim
prawdopodobieństwem pokryć się w skończonej
chwili z czasem przewidywalnym.
Czasy $S_k$ są czasami zatrzymania przez
adaptację, monotoniczność i prawostronną
ciągłość $N$, a $S_k>0$ p.n. dzięki $N_0=0$
i prawostronnej ciągłości. To kończy dowód.
\end{proof}

\begin{twierdzenie}[Warunkowe momenty skoku w czasie przewidywalnym]
\label{thm:predictable-jump-moments}
Niech $M$ będzie martyngałem c\`adl\`ag
z $M_0=0$ i $M_T\in L^2$, a $A=\langle M\rangle$
niech istnieje i spełnia $\E A_T<\infty$.
Niech $\tau$ będzie przewidywalnym czasem
zatrzymania, $0<\tau\le T$ p.n.
Definiujemy informację bezpośrednio przed $\tau$ przez
\[
\calF_{\tau-}=\sigma\bigl(\calF_0,\,
\{F\cap\{u<\tau\}:u\ge0,\ F\in\calF_u\}\bigr).
\]
Wówczas
\[
\E[\Delta M_\tau\mid\calF_{\tau-}]=0,\qquad
\E[(\Delta M_\tau)^2\mid\calF_{\tau-}]=\Delta A_\tau.
\]
\end{twierdzenie}
\begin{proof}
Niech $\tau_n$ będzie ciągiem zapowiadającym
$\tau$. Możemy przyjąć $0\le\tau_n<\tau\le T$
i $\tau_n\uparrow\tau$ poza jednym zbiorem
zerowym. Z dokładnością do uzupełnienia
\[
\calF_{\tau-}=\sigma\left(\bigcup_n\calF_{\tau_n}\right).
\]
Istotnie, dla $F\in\calF_{\tau_n}$ mamy
$F=\bigcup_{q\in\Q_{\ge0}}(F\cap\{\tau_n\le q\}
\cap\{q<\tau\})$ p.n., co daje jedną inkluzję.
Odwrotnie, dla $F\in\calF_u$ zdarzenie
$F\cap\{u<\tau\}$ jest sumą zbiorów
$F\cap\{u<\tau_n\}\in\calF_{\tau_n}$.
Ostatnia mierzalność wynika bezpośrednio
z testu definicji $\calF_{\tau_n}$.
Ponadto $\calF_0\subseteq\calF_{\tau_n}$.

Dla każdego przewidywalnego procesu $H$
zmienna $H_\tau$ jest $\calF_{\tau-}$-mierzalna.
Wystarczy sprawdzić to dla funkcji
charakterystycznych generatorów
$F\times(a,b]$, $F\in\calF_a$:
ich wartość w $\tau$ jest wskaźnikiem
$F\cap\{a<\tau\}\cap\{\tau\le b\}$.
Oba zdarzenia są $\calF_{\tau-}$-mierzalne.
Własność rozszerza się na całe $\sigma$-ciało
przewidywalne przez domknięcie na dopełnienia
i przeliczalne sumy. W szczególności
$A_\tau$ i $A_{\tau-}$ są mierzalne względem
$\calF_{\tau-}$ (proces $A_-$ jest przewidywalny),
podobnie $M_{\tau-}$.

Proces $M^2-A$ jest prawdziwym martyngałem
na $[0,T]$, ponieważ moduły jego wartości
w czasach zatrzymania są zmajoryzowane przez
$\sup_{s\le T}|M_s|^2+A_T\in L^1$
(twierdzenia~\ref{thm:doob-cont}
i~\ref{thm:local-mart-criteria}).
Dla $Y=M$ albo $Y=M^2-A$ ograniczone
próbkowanie opcjonalne daje
$\E[Y_\tau\mid\calF_{\tau_n}]=Y_{\tau_n}$.
Prawa strona zbiega p.n. i w $L^1$
do $Y_{\tau-}$ przez powyższą majorantę
(dla $Y=M$ wystarcza $\sup_{s\le T}|M_s|$).
Dla $F\in\bigcup_n\calF_{\tau_n}$
przechodzimy do granicy w równości
$\E[\chi_FY_\tau]=\E[\chi_FY_{\tau_n}]$
dla dostatecznie dużego $n$.
Twierdzenie o klasie monotonicznej rozszerza
równość na $\calF_{\tau-}$.
Stąd $\E[\Delta M_\tau\mid\calF_{\tau-}]=0$
i $\E[\Delta(M^2)_\tau-\Delta A_\tau
\mid\calF_{\tau-}]=0$.
Ponieważ
$\Delta(M^2)_\tau=2M_{\tau-}\Delta M_\tau+
(\Delta M_\tau)^2$,
a iloczyn jest całkowalny przez nierówność
Cauchy'ego--Schwarza, pierwsza równość
zeruje wyraz mieszany po warunkowaniu
(można najpierw obciąć $M_{\tau-}$).
Mierzalność $\Delta A_\tau$ daje drugą tezę.
\end{proof}

\begin{uwaga}[Intensywność nie musi być deterministyczna]
\label{rem:stochastic-counting-intensity}
Jeżeli kompensator procesu zliczającego ma postać
$A_t=\int_0^t a_s\,\mathrm d\lambda_1(s)$,
gdzie $a$ jest nieujemnym procesem przewidywalnym,
nazywamy $a$ intensywnością stochastyczną.
Tożsamość~\ref{thm:counting-compensation} daje
\[
\E\sum_{0<s\le T}H_s\Delta N_s
=\E\int_0^T H_s a_s\,\mathrm d\lambda_1(s).
\]
Intensywność jest jednoznaczna z dokładnością
do równości prawie wszędzie względem
$\lambda_1\otimes\Prob$: jednoznaczność
kompensatora daje równość wszystkich całek
do chwili $t$ na wspólnym zbiorze pełnej miary,
a więc równość gęstości prawie wszędzie
dla każdej takiej trajektorii i przez twierdzenie
Fubiniego. Nie wynika stąd niezależność
przyrostów procesu zliczającego.
\end{uwaga}

\begin{przyklad}[Losowa stała intensywność]
\label{ex:random-rate-counting}
Niech $\Lambda$ przyjmuje wartości $1,2$ z równymi
prawdopodobieństwami, a $P$ będzie niezależnym
procesem Poissona o intensywności $1$.
Definiujemy $N_t=P_{\Lambda t}$ i filtrację
generowaną przez $\Lambda$ oraz obserwacje
$N_s$ dla $s\le t$, z warunkami zwykłymi.
Warunkowo względem $\Lambda=\ell$ jest to
proces Poissona o intensywności $\ell$.
Niezależność przyrostów warunkowo względem
tej zmiennej daje
$\E[N_t-N_s\mid\calF_s]=\Lambda(t-s)$
oraz warunkową wariancję $\Lambda(t-s)$.
Równości zachowują się po zwykłym uzupełnieniu
filtracji: warunkujemy najpierw dla
$s_n\downarrow s$ i stosujemy odwrotne
warunkowe wartości oczekiwane oraz
prawostronną ciągłość procesów.
Stąd $N_t-\Lambda t$ i
$(N_t-\Lambda t)^2-\Lambda t$ są martyngałami.
Zatem kompensator $A_t=\Lambda t$ jest losowy.
Natomiast dla $h,k>0$, z prawa kowariancji
całkowitej,
\[
\operatorname{Cov}(N_h,N_{h+k}-N_h)
=hk\operatorname{Var}(\Lambda)=hk/4>0.
\]
Przyrosty na rozłącznych przedziałach
nie są więc niezależne bez warunkowania.
\end{przyklad}

\section{Złożony proces Poissona i jego fluktuacje}
\begin{twierdzenie}[Kompensowane sumy losowych skoków]
\label{thm:compound-poisson-brackets}
Niech $N$ będzie procesem Poissona o intensywności
$\lambda>0$, a $(Y_k)_{k\ge1}$ ciągiem niezależnych
zmiennych o jednakowym rozkładzie, niezależnym od $N$.
Załóżmy $\E Y_1^2<\infty$, oznaczmy
$\mu=\E Y_1$, $m_2=\E Y_1^2$ i zdefiniujmy
\[
J_t=\sum_{k=1}^{N_t}Y_k,\qquad
M_t=J_t-\lambda\mu t.
\]
Wtedy $M$ jest martyngałem w naturalnej filtracji
obserwacji zdarzeń i ich oznaczeń, z
\[
\E M_t^2=\lambda m_2t,\qquad
[M]_t=\sum_{k=1}^{N_t}Y_k^2,\qquad
\langle M\rangle_t=\lambda m_2t.
\]
\end{twierdzenie}
\begin{proof}
Stosujemy własność niezależnych przyrostów
w filtracji obserwowanych oznaczeń z rozdziału
\ref{ch:poisson-process}. W szczególności nie
ujawniamy z góry całego ciągu $(Y_k)$.
Przyrost $J_t-J_s$ jest niezależny od informacji
do chwili $s$. Warunkowo względem $N_t-N_s=n$ jest sumą
$n$ niezależnych kopii $Y_1$. Zatem
$\E(J_t-J_s)=\lambda(t-s)\mu$ i $M$ jest martyngałem.
Z prawa wariancji całkowitej
\[
\begin{aligned}
\Var(J_t-J_s)
&=\E[(N_t-N_s)\Var(Y_1)]
 +\Var((N_t-N_s)\mu)\\
&=\lambda(t-s)(\Var Y_1+\mu^2)
=\lambda(t-s)m_2.
\end{aligned}
\]
Stąd $\E[M_t^2\mid\calF_s]
=M_s^2+\lambda m_2(t-s)$, a więc
$M_t^2-\lambda m_2t$ jest martyngałem.
Trajektoria $M$ ma pomiędzy skokami wariację skończoną,
a jej przyrosty w chwilach zdarzeń są równe
kolejnym wartościom $Y_k$ (oznaczenie zerowe
nie daje niezerowego skoku).
Z lematu~\ref{lem:cadlag-finite-variation-qv}
mamy $[M]_t=\sum_{k\le N_t}Y_k^2$.
Stosując ten sam rachunek warunkowej średniej
do oznaczeń $Y_k^2$ (wystarcza ich całkowalność),
otrzymujemy martyngał $[M]_t-\lambda m_2t$.
Stąd również $M^2-[M]$ jest martyngałem.
Deterministyczny proces $\lambda m_2t$ jest
ciągły, niemalejący i przewidywalny, więc spełnia
definicję przewidywalnego nawiasu $M$.
\end{proof}

\begin{twierdzenie}[Tożsamość kompensacyjna dla oznaczeń]
\label{thm:marked-poisson-compensation}
Niech $S_k$ będą czasami zdarzeń procesu Poissona
o intensywności $\lambda$, a $Y_k$ niezależnymi
oznaczeniami o wspólnym rozkładzie $\nu$ na $\R$,
niezależnymi od procesu Poissona. Filtracja
ujawnia oznaczenie $Y_k$ dopiero w $S_k$.
Dla nieujemnej funkcji $H(t,\omega,y)$
mierzalnej względem iloczynu $\sigma$-ciała
przewidywalnego i $\mathcal B(\R)$ zachodzi
\[
\E\sum_{k:S_k\le T}H(S_k,\omega,Y_k)
=\lambda\E\int_0^T\int_{\R}H(t,\omega,y)
\,\mathrm d\nu(y)\,\mathrm d\lambda_1(t).
\]
\end{twierdzenie}
\begin{proof}
Najpierw bierzemy
$H(t,\omega,y)=\chi_F(\omega)\chi_{(a,b]}(t)
\chi_B(y)$, $F\in\calF_a$, $0\le a<b\le T$,
$B\in\mathcal B(\R)$.
Z własności oznaczonych przyrostów
(rozdział~\ref{ch:poisson-process})
warunkowa średnia liczby przyszłych zdarzeń
z oznaczeniem w $B$ wynosi
$\lambda(b-a)\nu(B)$: warunkowo względem
liczby przyszłych zdarzeń oznaczenia są
niezależnymi kopiami o prawie $\nu$.
Po pomnożeniu przez $\chi_F$ i wzięciu
wartości oczekiwanej otrzymujemy tezę.
Te prostokąty tworzą układ zamknięty na
skończone przecięcia, generujący rozważane
$\sigma$-ciało dla $0<t\le T$.
Obie strony dla funkcji charakterystycznych
są skończonymi miarami (ich całkowita masa
wynosi $\lambda T$); lewa strona jest
mierzalna najpierw dla generatorów,
a następnie przez klasę monotoniczną.
Z jednoznaczności miar równość zachodzi
dla wszystkich zbiorów mierzalnych.
Aproksymacja funkcjami prostymi i zbieżność
monotoniczna dają wynik.
\end{proof}

\begin{zadanie}[Ryzyko portfela ze skokami]
Niech $N$ będzie procesem Poissona o intensywności $2$,
a skoki $Y_k$ przyjmują wartości $+3$ i $-1$,
każdą z prawdopodobieństwem $1/2$.
Wyznaczyć wartość oczekiwaną i wariancję $J_t$,
a także $\E[J_t\mid N_t=3]$.

\emph{Rozwiązanie.} Mamy
$\mu=(3-1)/2=1$, $m_2=(9+1)/2=5$.
Stąd $\E J_t=2t$, $\Var(J_t)=2t\cdot5=10t$.
Przy warunku $N_t=3$ suma zawiera dokładnie
trzy niezależne oznaczenia, zatem
$\E[J_t\mid N_t=3]=3\mu=3$.
To nie jest to samo co $\E J_t=2t$:
warunkowanie ujawnia informację o liczbie skoków.
\end{zadanie}


\begin{zadanie}[Rzeczywista i przewidywalna wariacja --- dwa różne wyniki]
Niech $N$ będzie procesem Poissona o intensywności $2$,
a $M_t=N_t-2t$. Obliczyć prawdopodobieństwo zdarzenia
$\{[M]_1=0\}$, wartości oczekiwane $[M]_1$ oraz
$\langle M\rangle_1$, i wyjaśnić, dlaczego na tym
zdarzeniu oba nawiasy są różne.

\emph{Rozwiązanie.} Ponieważ $[M]_1=N_1$,
$\Prob([M]_1=0)=\Prob(N_1=0)=e^{-2}$.
Natomiast $\langle M\rangle_1=2$ deterministycznie,
także gdy do chwili $1$ nie wystąpił żaden skok.
Oba nawiasy mają jednak tę samą wartość oczekiwaną:
$\E\bigl([M]_1\bigr)=\E N_1=2
=\E\bigl(\langle M\rangle_1\bigr)$.
Natomiast $\E M_1=0$; są to trzy różne wielkości.
Jest to najprostsza ilustracja różnicy między wielkością
zrealizowaną a jej przewidywalną kompensacją.
\end{zadanie}


\section{Skończone sumy przewidywalne dla procesów ze skokami}
\begin{twierdzenie}[Izometria kompensowanej sumy]
\label{thm:compound-poisson-isometry}
Niech $M$ będzie kompensowanym procesem z twierdzenia
\ref{thm:compound-poisson-brackets}, a $0=t_0<\cdots<t_n=T$.
Jeśli $H_i\in L^\infty(\calF_{t_i})$, to
$S=\sum_{i=0}^{n-1}H_i(M_{t_{i+1}}-M_{t_i})$
spełnia
\[
\E S=0,\qquad
\E S^2=\lambda m_2\sum_{i=0}^{n-1}
(t_{i+1}-t_i)\E H_i^2.
\]
\end{twierdzenie}
\begin{proof}
Pierwsza równość i zanik wyrazów mieszanych w
$\E S^2$ wynikają z niezależnych wycentrowanych
przyrostów, dokładnie jak w lemacie o skończonych
sumach przewidywalnych z rozdziału
\ref{ch:quadratic-cont}. Dla każdego $i$
\[
\E[H_i^2(M_{t_{i+1}}-M_{t_i})^2]
=\E H_i^2\,\lambda m_2(t_{i+1}-t_i),
\]
ponieważ kwadrat przyrostu jest niezależny od
$\calF_{t_i}$. Po zsumowaniu dostajemy tezę.
\end{proof}

\begin{twierdzenie}[Izometria dla kompensowanego procesu zliczającego]
\label{thm:counting-compensated-isometry}
Przy założeniach twierdzenia
\ref{thm:counting-compensation} załóżmy dodatkowo,
że $A$ jest ciągły oraz $M=N-A$ jest martyngałem
kwadratowo całkowalnym na $[0,T]$.
Wtedy
\[
[M]_t=N_t,\qquad \langle M\rangle_t=A_t.
\]
Dla $0=t_0<\cdots<t_m=T$ i ograniczonych
$\calF_{t_i}$-mierzalnych $H_i$ zachodzi
\[
\E\left(\sum_i H_i(M_{t_{i+1}}-M_{t_i})\right)^2
=\E\sum_i H_i^2(A_{t_{i+1}}-A_{t_i}).
\]
\end{twierdzenie}
\begin{proof}
Z lematu~\ref{lem:cadlag-finite-variation-qv}
mamy $[M]=N$, gdyż $M$ ma lokalnie skończoną
wariację, a ciągłość $A$ daje $\Delta M=\Delta N$.
Udowodnimy bez całki It\^o, że $M^2-N$ jest
lokalnym martyngałem. Dla trajektorii o skończonej
wariacji zachodzi tożsamość
\[
M_t^2=2\int_{(0,t]}M_{s-}\,\mathrm dM_s+
\sum_{0<s\le t}(\Delta M_s)^2 .
\]
Wynika ona z teleskopowania kwadratów na
podziałach: sumy kwadratów zbiegają do $[M]_t$,
a sumy z lewym końcem przyrostu zbiegają
do całki względem miary znakowanej $\mathrm dM$.
Istotnie, na $(t_i,t_{i+1}]$ funkcja
podcałkowa jest równa $M_{t_i}$; dla każdego
$s>0$ zbiega ona do $M_{s-}$ przy zanikającej
średnicy, gdyż $t_i<s$. Zbieżność zmajoryzowana
względem skończonej miary wariacji daje tezę,
także dla podziału kończącego się w $t$.

Niech $\rho_k=\inf\{t:N_t+A_t\ge k\}\wedge k$.
Proces $N+A$ jest niemalejący i adaptowany,
jego skoki mają rozmiar jeden, zatem
$(N+A)_{t\wedge\rho_k}\le k+1$.
Czasy $\rho_k\uparrow\infty$ p.n.
Zatrzymany $M^{\rho_k}$ jest ograniczonym
martyngałem i ma wariację nie większą niż $k+1$.
Dla $0\le s<t$ i $F\in\calF_s$ stosujemy
lemat~\ref{lem:predictable-fv-mean-zero}
do $D=M^{\rho_k}$ i ograniczonego
procesu przewidywalnego
$\chi_F\chi_{(s,t]}D_-$.
Otrzymujemy, że trajektoryjna całka
$\int_{(0,t]}D_{u-}\,\mathrm dD_u$
jest martyngałem (adaptacja wynika też
z granic skończonych sum).
Tożsamość kwadratowa daje martyngał
$(M^{\rho_k})^2-N^{\rho_k}$.
Zatem $M^2-N$ jest lokalnym martyngałem,
a dodając $N-A$ otrzymujemy lokalny
martyngał $M^2-A$.

Na $[0,T]$ rodzina wartości $M^2-A$ w
ograniczonych czasach zatrzymania jest
jednostajnie całkowalna, ponieważ jej moduły
są ograniczone przez
$\sup_{s\le T}|M_s|^2+A_T\in L^1$
(nierówność Dooba, twierdzenie
\ref{thm:doob-cont}). Jest więc prawdziwym
martyngałem na $[0,T]$ na mocy kryterium
lokalnego martyngału z rozdziału
\ref{ch:mart-cont}.
Proces $A$ spełnia definicję nawiasu
przewidywalnego, a jednoznaczność daje
wniosek~\ref{cor:predictable-compensator-unique}.
Rozwijając kwadrat przyrostu i warunkując
względem jego początku otrzymujemy
$\E[(M_v-M_u)^2\mid\calF_u]
=\E[A_v-A_u\mid\calF_u]$.
Ortogonalność różnych przyrostów
z przewidywalnymi współczynnikami
oraz własność wieży dają żądaną izometrię.
Nie potrzeba niezależności przyrostów $N$.
\end{proof}

\section{Wykładnik procesu Poissona i oszacowania ogonów}
\begin{twierdzenie}[Martyngał wykładniczy procesu Poissona]
\label{thm:poisson-exponential-mart}
Niech $N$ będzie procesem Poissona o intensywności
$\lambda>0$ w uzupełnionej filtracji naturalnej.
Dla każdej $\theta\in\R$ proces
\[
Z_t(\theta)=\exp\{\theta N_t-\lambda t(e^\theta-1)\}
\]
jest dodatnim martyngałem o wartości oczekiwanej $1$.
Dla $T>0$ oraz $a>\lambda T$ zachodzi
\[
\Prob(N_T\ge a)\le
\exp\left(-a\log\frac{a}{\lambda T}+a-\lambda T\right).
\]
\end{twierdzenie}
\begin{proof}
Dla $s<t$ niezależność przyrostu
$U=N_t-N_s\sim\mathrm{Poisson}(\lambda(t-s))$
i jego funkcja tworząca
$\E e^{\theta U}=\exp(\lambda(t-s)(e^\theta-1))$
dają $\E[Z_t(\theta)\mid\calF_s]=Z_s(\theta)$.
Dla $\theta>0$ nierówność Markowa zastosowana do
$e^{\theta N_T}$ prowadzi do
$\Prob(N_T\ge a)\le
\exp(-\theta a+\lambda T(e^\theta-1))$.
Minimalizujemy wykładnik względem $\theta>0$:
pochodna wynosi $-a+\lambda Te^\theta$, więc
minimum przypada przy $\theta=\log(a/(\lambda T))$.
Podstawienie daje tezę. Jest to ścisły przykład tego,
jak martyngał wykładniczy prowadzi do nierówności
koncentracyjnej, a nie tylko do wzorów na momenty.
\end{proof}

\begin{twierdzenie}[Martyngał wykładniczy przy losowym kompensatorze]
\label{thm:counting-exponential-bound}
Niech $N,A$ spełniają założenia twierdzenia
\ref{thm:counting-compensation}, a $A$ będzie
ciągły. Dla $\theta\in\R$ proces
\[
Z_t=\exp\{\theta N_t-(e^\theta-1)A_t\}
\]
jest dodatnim lokalnym martyngałem, zatem
nadmartyngałem o $\E Z_t\le1$.
Dla $M=N-A$, $x,v>0$ i $T>0$ zachodzi
\[
\Prob\left(\sup_{t\le T\wedge\eta_v}M_t\ge x\right)
\le \exp\{-v h(x/v)\},
\qquad \eta_v=\inf\{t:A_t>v\},
\qquad h(r)=(1+r)\log(1+r)-r .
\]
\end{twierdzenie}
\begin{proof}
Trajektorie $Z$ mają lokalnie skończoną
wariację, ponieważ $N$ i $A$ mają tę własność.
Na odcinkach bez zdarzeń $N$ reguła łańcuchowa
dla ciągłej funkcji niemalejącej $A$ daje
$\mathrm dZ=-(e^\theta-1)Z\,\mathrm dA$.
Reguła ta nie wymaga różniczkowalności $A$:
teleskopujemy funkcję wykładniczą po
podziałach i szacujemy błąd liniowego
przybliżenia przez stałą razy
$\max_i|\Delta_i A|\sum_i|\Delta_i A|$,
który dąży do zera przez ciągłość $A$.
W chwili skoku $\Delta N=1$ mamy
$\Delta Z=(e^\theta-1)Z_-$.
Miara $\mathrm dA$ nie ma atomów,
więc $Z$ i $Z_-$ są równe względem tej miary
prawie wszędzie. Sumując części ciągłe
i skoki otrzymujemy trajektoryjnie
\[
Z_t=1+(e^\theta-1)\int_{(0,t]}Z_{s-}\,\mathrm d(N-A)_s .
\]
Nie używamy tutaj całki It\^o.
Lokalizujemy czasami
$\rho_k=k\wedge\inf\{t:N_t+A_t\ge k\}$.
Procesy $(N-A)^{\rho_k}$ są martyngałami
o wariacji ograniczonej przez $k+1$.
Wartości $Z_{s-}$ przed zatrzymaniem
są ograniczone przez stałą zależną od $k,\theta$.
Zatrzymana tożsamość jest zatem całką
względem $(N-A)^{\rho_k}$ z ograniczonym
przewidywalnym współczynnikiem
$(Z^{\rho_k})_-$.
Lemat~\ref{lem:predictable-fv-mean-zero},
stosowany także po pomnożeniu współczynnika
przez $\chi_F\chi_{(u,t]}$, $F\in\calF_u$,
daje własność martyngałową $Z^{\rho_k}$.
Ponieważ $\rho_k\uparrow\infty$ p.n.,
$Z$ jest lokalnym martyngałem.
Nieujemny lokalny martyngał jest nadmartyngałem
(twierdzenie~\ref{thm:positive-local-supermart}).

Dla $\theta>0$ zapisujemy
$Z_t=\exp\{\theta M_t-(e^\theta-1-\theta)A_t\}$.
Czas $\eta_v$ jest czasem zatrzymania dzięki
adaptacji, ciągłości $A$ i prawostronnej
ciągłości filtracji, jak w dowodzie
twierdzenia~\ref{thm:dambis-dubins-schwarz}.
Dla $t\le T\wedge\eta_v$ mamy $A_t\le v$.
Biorąc supremum, na zdarzeniu z tezy mamy
$\sup_{t\le T} Z_t\ge
e^{\theta x-(e^\theta-1-\theta)v}$
także wtedy, gdy supremum $M$ nie jest osiągnięte.
Nierówność Ville'a
(twierdzenie~\ref{thm:ville-cont}) daje
oszacowanie
\[
\exp\{-\theta x+(e^\theta-1-\theta)v\}.
\]
Wybór $\theta=\log(1+x/v)>0$ minimalizuje
wykładnik i daje $-vh(x/v)$.
Supremum procesu zatrzymanego $M^{\eta_v}$
jest mierzalne, ponieważ można je obliczyć
na przeliczalnym gęstym zbiorze chwil
wraz z końcem $T$, dzięki prawostronnej
ciągłości.
\end{proof}

\begin{uwaga}[Znaczenie ograniczenia nawiasu]
W przypadku Poissona $A_t=\lambda t$,
więc przy $v=\lambda T$ ostatnia nierówność
dotyczy maksimum skompensowanego procesu
na całym $[0,T]$, a nie tylko jego końcowej
wartości. Dla losowego $A$ oszacowanie
kontroluje przekroczenia przed wyczerpaniem
budżetu kompensatora $v$.
Nie zakładamy, że dodatni lokalny martyngał
wykładniczy jest zawsze prawdziwym martyngałem.
\end{uwaga}

\section{Martyngał ciągły i skokowy w jednym modelu}

\begin{twierdzenie}[Dodanie skończonej wariacji do procesu ciągłego]
\label{thm:finite-variation-jumps-bracket}
Niech $C$ będzie procesem o ciągłych trajektoriach, dla którego
wariacja kwadratowa $[C]$ istnieje w sensie zbieżności ucp.
Niech $A$ będzie procesem c\`adl\`ag o skończonej wariacji
całkowitej na każdym przedziale $[0,T]$, prawie na pewno.
Liczba skoków może być nieskończona. Wtedy $X=C+A$ ma wariację
kwadratową i zachodzi, z dokładnością do nierozróżnialności,
\[
[X]_t=[C]_t+\sum_{0<s\le t}(\Delta A_s)^2,
\qquad t\ge0.
\]
Nie zakładamy niezależności $C$ i $A$.
\end{twierdzenie}
\begin{proof}
Ustalmy $T>0$ i podział $\pi$. Dla każdego $t\le T$
rozwinięcie kwadratu daje
\[
V_t^\pi(C+A)=V_t^\pi(C)+V_t^\pi(A)
 +2\sum_i\Delta_i^tC\,\Delta_i^tA,
\]
gdzie $\Delta_i^tC=C_{t\wedge t_{i+1}}-C_{t\wedge t_i}$
i analogicznie definiujemy $\Delta_i^tA$. Dla każdej
ustalonej trajektorii $A$ o skończonej wariacji całkowitej
\[
\sup_{t\le T}\left|\sum_i\Delta_i^tC\,\Delta_i^tA\right|
\le \omega_C(|\pi|)\operatorname{Var}_{[0,T]}(A),
\qquad
\omega_C(\delta):=
\sup_{\substack{u,v\in[0,T]\\|u-v|\le\delta}}|C_u-C_v|.
\]
Ponieważ $C$ ma ciągłe trajektorie, $\omega_C(\delta)\to0$
przy $\delta\downarrow0$. Iloczyn po prawej stronie zbiega
więc do zera prawie na pewno, jednostajnie po $t$.
Z lematu~\ref{lem:cadlag-finite-variation-qv}
$V^\pi(A)$ zbiega jednostajnie
prawie na pewno do sumy kwadratów skoków, natomiast
$V^\pi(C)\to[C]$ w ucp z założenia.
Sumując te trzy zbieżności, otrzymujemy tezę.
Ważne jest, że ciągłość $C$ wyklucza wspólny skok;
w przeciwnym razie pojawiłyby się również iloczyny
skoków $2\Delta C_s\Delta A_s$.
\end{proof}

\begin{stwierdzenie}[Brak kowariacji niezależnych składników]
Niech $W$ będzie procesem Wienera, a $N$ niezależnym
od niego procesem Poissona o intensywności $\lambda>0$.
Dla $M_t=N_t-\lambda t$ i dowolnych deterministycznych
ciągów podziałów $\pi$ przedziału $[0,T]$
o średnicach dążących do zera zachodzi
\[
\sum_i(W_{t_{i+1}}-W_{t_i})
(M_{t_{i+1}}-M_{t_i})\longrightarrow0
\quad\text{w }L^2.
\]
W szczególności wariacja kwadratowa sumy
$X_t=W_t+M_t$ ma postać $[X]_t=t+N_t$.
Sama tożsamość wariacji kwadratowej nie wymaga
niezależności $W$ i $N$: wynika także z ciągłości $W$
i skończonej wariacji $M$ na przedziałach zwartych.
\end{stwierdzenie}
\begin{proof}
Niezależne przyrosty obu procesów oraz niezależność
całych procesów dają zanik mieszanych momentów
między różnymi komórkami siatki. Dla pojedynczej komórki
długości $h_i$ mamy
$\E(\Delta_iW)^2=h_i$ oraz
$\E(\Delta_iM)^2=\lambda h_i$.
Zatem
\[
\E\left|\sum_i\Delta_iW\Delta_iM\right|^2
=\lambda\sum_i h_i^2\le\lambda T|\pi|\longrightarrow0.
\]
Z elementarnej tożsamości
$\sum(\Delta X)^2=\sum(\Delta W)^2+
2\sum\Delta W\Delta M+\sum(\Delta M)^2$
oraz wyników wcześniejszych rozdziałów wynika
$[X]_t=t+N_t$. Dla pełnej ścisłości zbieżność ucp
wynika także bezpośrednio z twierdzenia
\ref{thm:finite-variation-jumps-bracket}: trajektoria
$M=N-\lambda t$ ma skończoną wariację i skończenie
wiele skoków na $[0,T]$, zaś $W$ jest ciągły.
Oszacowanie $L^2$ powyżej daje ponadto ilościową kontrolę
mieszanych sum w końcowej chwili $T$.
To rozdziela dwa fakty: zanik kowariacji jest tutaj
własnością regularności trajektorii, natomiast przedstawione oszacowanie $L^2$
korzysta z niezależności.
Jeżeli $W$ i $M$ są martyngałami względem tej samej
filtracji i są kwadratowo całkowalne, to
$WM$ również jest martyngałem w przypadku niezależnych
procesów w ich wspólnej uzupełnionej filtracji
naturalnej: warunkowanie względem wspólnej przeszłości
redukuje się do warunkowania w obu niezależnych
filtracjach (najpierw na prostokątach zdarzeń).
Wtedy $X^2-(1+\lambda)t$ jest martyngałem przez
rozwinięcie kwadratu i twierdzenie
\ref{thm:poisson-brackets}, a zatem
$\langle X\rangle_t=(1+\lambda)t$.
\end{proof}

\section{Co przenosimy do tomu III?}
Tom II został zbudowany wokół pytania \emph{jak zachowują
się procesy losowe}. Badaliśmy ich rozkłady, trajektorie,
filtracje, czasy zatrzymania, własność Markowa, strukturę
odnowy i rozgałęziania oraz martyngały i wariację
kwadratową. Kolejny tom zmieni pytanie na
\emph{jak przeprowadzać analizę na takich trajektoriach}.

\begin{uwaga}[Granica pomiędzy procesami a analizą stochastyczną]
W tomie III od podstaw skonstruujemy całkę względem
procesu Wienera dla procesów przewidywalnych, wzór
Itô, całkę Stratonovicha, lokalne czasy i formułę
Tanaki, a następnie całkowanie względem ogólnych
martyngałów i półmartyngałów oraz równania stochastyczne.
Później powrócimy do miary Wienera, całek funkcyjnych,
wzoru Feynmana--Kaca i relacji do formalnych całek
Feynmana po trajektoriach. Skończone sumy przewidywalne
z rozdziałów 13--14 są zapowiedzią tej konstrukcji,
ale same nie stanowią jeszcze ogólnej teorii całki Itô.
\end{uwaga}

\begin{zadanie}[Dwie różne miary fluktuacji]
Niech $W$ będzie procesem Wienera, $N$ niezależnym
procesem Poissona o intensywności $3$ oraz
$X_t=W_t+(N_t-3t)$. Obliczyć $[X]_t$ i średnią
$\E[X_t^2]$.

\emph{Rozwiązanie.} Kwadratowa zmienność ciągłej części
wynosi $t$, a skoki części poissonowskiej mają rozmiar
jeden, więc $[X]_t=t+N_t$. Z niezależności
$W_t$ oraz $N_t-3t$ mamy
$\E X_t=0$ i
$\E X_t^2=\Var(W_t)+\Var(N_t)=t+3t=4t$.
Także $\E[X]_t=t+\E N_t=4t$.
Zauważamy jednak, że $[X]_t$ jest losowy, podczas
gdy $\langle X\rangle_t=4t$ jest deterministyczny.
Tempo wzrostu przewidywalnego nawiasu wynosi $4$.
\end{zadanie}

